\textcolor{black}{\section{Introducción}}
\noindent En sistemas electrónicos de control y automatización, es común utilizar dispositivos que permitan manejar cargas de mayor potencia mediante señales de bajo nivel, como las provenientes de microcontroladores o circuitos lógicos. En este contexto, los relevadores constituyen uno de los elementos más empleados debido a su capacidad para establecer o interrumpir conexiones eléctricas mediante un mecanismo electromagnético. Un relevador, también llamado relé, es un interruptor controlado eléctricamente diseñado para abrir o cerrar un circuito cuando circula corriente por una bobina interna, lo que provoca un efecto de atracción magnética que desplaza un contacto mecánico \textbf{[3]}. Su funcionamiento se basa en la inducción electromagnética: cuando el devanado recibe corriente, se genera un campo magnético que mueve una armadura metálica, permitiendo conmutar entre estados de encendido y apagado de la carga conectada \textbf{[4]}. Este comportamiento lo convierte en un dispositivo ideal para aislar eléctricamente un circuito de control de un circuito de potencia.

La importancia de los relevadores radica en que permiten controlar cargas de corriente alterna o continua con niveles de potencia elevados, utilizando señales de mando de baja energía. Esto los hace indispensables en aplicaciones donde se requiere protección, aislamiento o control remoto de dispositivos eléctricos, como motores, lámparas, sistemas de seguridad, automatización industrial y domótica \textbf{[1]}. Entre sus aplicaciones más comunes se encuentran: sistemas de iluminación automatizada, control de bombas de agua, sistemas de alarma, arranque de motores eléctricos y automatización residencial. Además, existen diversos tipos de relevadores, entre los que destacan los electromecánicos, de estado sólido (SSR), térmicos, de láminas (reed) y relés temporizados, cada uno diseñado para condiciones específicas de carga, velocidad de respuesta y durabilidad \textbf{[2]}.

En los sistemas electrónicos de control, es frecuente la necesidad de activar cargas eléctricas de potencia como lámparas, motores o alarmas, utilizando señales de bajo nivel provenientes de circuitos lógicos o microcontroladores. Sin embargo, dichos dispositivos digitales no pueden suministrar corrientes suficientemente altas para activar directamente estos actuadores. Por ello, el transistor BJT (Bipolar Junction Transistor) se emplea como intermediario para amplificar corriente y funcionar como conmutador electrónico. Su función consiste en permitir o bloquear el paso de corriente desde el colector hacia el emisor, dependiendo de la corriente que circule por la base \textbf{[4]}.

Cuando se utiliza como interruptor, el transistor trabaja únicamente en dos regiones: corte (no conduce, equivalente a un interruptor abierto) y saturación (conduce totalmente, simulando un interruptor cerrado). Estas propiedades permiten generar conmutaciones precisas y rápidas, lo que lo hace ideal para el control de relés \textbf{[1]}. El relé, por su parte, es un interruptor de operación electromagnética, compuesto por una bobina, contactos móviles y un resorte. Cuando la bobina recibe corriente, genera un campo magnético que mueve los contactos y permite abrir o cerrar un circuito de potencia. Esto permite activar cargas de corriente alterna (CA) mediante señales de corriente directa (CD), manteniendo aislamiento eléctrico entre la etapa de control y la de potencia \textbf{[2]}.

Los relés son fundamentales en automatización, seguridad, control industrial, robótica, alarmas contra incendio, sistemas HVAC y domótica, ya que permiten la operación remota y segura de cargas de potencia \textbf{[3]}. Hoy en día, los relés pueden controlarse mediante plataformas programables como Arduino, combinando sensores como el ultrasónico o el LM35 para desarrollar aplicaciones inteligentes como sistemas de iluminación automática, alarma por temperatura, detección de presencia y control por temporización. Esto demuestra cómo el transistor, el relé y el microcontrolador conforman una arquitectura híbrida entre electrónica analógica y digital.

Es importante tener en consideración como el transistor puede ser controlado por señales digitales provenientes de una placa Arduino, permitiendo implementar sistemas temporizados, sensores de ultrasonido para detección de proximidad y sensores LM35 para monitoreo de temperatura, lo que demuestra su amplia versatilidad en aplicaciones de automatización. Como complemento, es posible notar el comportamiento dinámico del circuito mediante mediciones de corriente de base (IB), corriente de colector (IC) y los voltajes VBE y VCE, tanto en simulación como en laboratorio.

